\documentclass[conference]{IEEEtran}

\begin{document}

\title{Group 7 Project Proposal}

\author{
    \IEEEauthorblockN{Mohammad Islam, Macy Narang, Ken Jiang, and Jinxin Lin}
    \IEEEauthorblockA{
        \textit{Computer Science} \\
        \textit{CUNY Hunter College} \\
        mohammad.islam47@login.cuny.edu,\quad macy.narang99@login.cuny.edu \\
        ken.jiang24@login.cuny.edu,\quad jinxin.lin07@login.cuny.edu
    }
}

\maketitle

\section{Introduction}

Several methods for estimating the win probability of a specific NFL team given
only the current state of an ongoing game are already widely used in sports
analytics. However, such calculations of win probability are not informative of
how that probability could change across the range of possible next plays and
their outcomes, nor does the win probability update until after a play has
already been completed. Thus, the aim of this project is to address this
problem by estimating a distribution of win probabilities associated with a
current game state conditional on plausible next plays before the next play has
even finished. This service will be delivered through a web tool that can be
used by general audiences, professional analysts, and NFL personnel alike.

\section{Intended Outcome}

If the models and web tool work as intended, they will give users insight into
how win probability changes throughout a game depending on what plays are
called. The information generated by our tool can be used to assist football
coaches in their play-making decisions, for educational purposes among fans and
new coaching personnel, and to satisfy general curiosity among users. In
addition, the fact that our tool will enable users to be proactive about game
information rather than reactive to play outcomes means that it could also
narrow the gap between professional traders---including hedge funds---and
regular people who participate in sports-related prediction market trading.
Currently, most professional traders and hedge funds buy access to the lowest
latency possible real-time data or have people attend games for them and report
on play outcomes. Their intent in doing so is to obtain and react to the result
of a completed play before their competition can, meaning that they essentially
dominate the market at present. However, since our web tool allows people to
make predictions before a play is finished, it can give less sophisticated
traders a fighting chance against people who are paying for low-latency
real-time win probability that updates after each play. For example, if a
prediction market trader has already made an input to our web tool before a
play has been completed, and if they can then recognize what type of play is
unfolding in real time, they could change how they are trading or initiate a
trade in response to the win probability displayed on our website that is
associated with that type of play. Since most football plays last for several
seconds, and since even the lowest-latency post-play updates are delayed to
some extent, it should very often be feasible for our users to make inputs to
our model and see its predictions before traders who react to post-play win
probability have even made their trades, no matter how fast the latter party's
data is obtained. This would contribute to leveling the playing field in sports
trading. In sum, if we can make a tool that promotes fairness in markets and
encourages learning and experimentation across a relatively varied user base,
we will have achieved our goal in undertaking this project.

\section{Approach}

The models that calculate the probabilities mentioned above will be trained on
game-state data---including down, distance, time remaining, etc.---obtained
from a Python port of the nflverse called nflreadpy. The nflverse offers a
dataset of historical NFL play-by-play data up to and including this season.
This data set will be converted to and accessed using pandas dataframes for
initial testing, but since it contains several tens of thousands of rows, the
final product will instead use polars for increased efficiency. The most
essential feature in our model, which is the next play type, is not a column in
this data set. However, since each row is organized sequentially and contains
the current play, we can engineer that feature ourselves. Our dependent
variable will be the outcome of the game, specifically whether or not the home
team won, but the outputs will be the probabilities associated with the home
team winning conditional on the next play.

We will build our models in Python using multiple different structures,
including neural networks, logistic regressions, and possibly k-nearest
neighbors and decision trees. The final model(s) should take vectors of live
game-state information as input, and deliver estimated win probabilities
corresponding to plausible play types as output. The model(s) will be
integrated into an interactive web tool that will allow users to input to the
model by selecting an ongoing game and importing the live state of that game
through web scraping or APIs, or by manually entering information in fields
that correspond to the input variables of the model. Upon receiving user input,
the webpage will display the resulting win probabilities conditioned on the
next play. If we end up picking multiple final models, we will put the win
probabilities associated with each model on separate lines and label them
accordingly.

During the testing phase, we will evaluate the models' performance by
separating the data into training and test sets, and evaluating the performance
of the models on the test set using metrics that are designed to calculate loss
on probability estimates, such as Brier score loss. We will also perform
real-world sanity testing by watching ongoing games, making inputs to the
model, and seeing whether the outputs align with ESPN's post-play win
probability predictions and our expectations. We could perform a similar
progress on historical games as well. Finally, if allowed, we will have outside
parties with varying levels of football knowledge test our product and give us
feedback on how intuitive it is to use.

To further aid users in usage and understanding of our win probability
estimates, we will display information and statistics about the players on the
offensive team that are relevant to the most successful play types, recent
head-to-head matchup data between the teams, and team rankings and statistics,
among other things. This data can be obtained using datasets, APIs, or web
scraping.

\subsection{Roles}

Macy and Jinxin's roles in this project will primarily be to work with the
models in Python. Macy and Jinxin will each pick at least one type of model to
experiment with; for example, Macy might focus on logistic regressions and
k-nearest neighbors while Jinxin works on neural networks and decision trees.
We will build our models using scikit-learn for traditional machine learning
algorithms and TensorFlow's Keras API for neural networks. We may research and
incorporate other python packages if necessary for data processing or model
construction. In the early stages, Macy and Jinxin will each build a
preliminary model that will be used by the other group members for testing and
development in their own roles while researching different kinds of model
architecture and machine learning algorithms. Then, Macy will collaborate with
Ken to select the features that will be inputted to the model, and Jinxin will
research how to web scrape or use an API to get live game state data. After
this, they will each construct several models, run simulations on hypothetical
game states, perform both real-world and mathematical evaluations of the
models, and give each other feedback. Based on their results, they will make a
joint decision on which model(s) should be used in the final product.

Ken's role will be to handle the backend. He will use Python packages to build
and manage the data pipeline. He will also load and preprocess the dataset in
Polars, assist in choosing the model features, and construct a the final
dataframe. The final dataframe that the model will be trained on will contain
only the variables needed for the model, all but two of which---next play and
game outcome---are also the ones that will be obtained from APIs or inputted by
users. This smaller dataset will be stored in a PostgreSQL database and hosted
on Supabase using SQL. Finally, Ken will use FastAPI to provide API endpoints
that receive the game state data from the frontend, pass it to the Python
model(s), and return the results to the frontend as JSON.

Mohammad will manage the frontend and design the user interface of our web tool
to align with our project goals and feel intuitive and engaging to potential
users. The user interface, functionality, and features will be updated over
time in a way that adapts to the progress that is made by other group members.
Additionally, Mohammad will handle frontend scripting, hosting the website and
selecting a platform on which to do so, integrating APIs into the website, and
securing API keys that should not be accessible to regular users. He will use
React, Vite, HTML, CSS, and TypeScript to develop the web tool.

If any of the group members has time to take on additional roles, they can work
on backtesting, finding APIs and writing web scripts to integrate into the
website that can provide information about specific current games and relevant
players, or refining or changing our method of obtaining real-time game data.

\subsection{Timeline}

By mid-October, we are aiming to have finished preprocessing, selecting model
features, creating a working webpage, and building the data pipeline for the
model. In early November, we will ideally have finished building several
possible models, updated the frontend of the site to reflect our backend and
modeling progress, and made sure the dataset and model code are properly
integrated into the web tool. After this, we will undergo the testing
procedures mentioned in previous sections with the aim of finalizing our
webpage and model by the last week of November. Throughout the project, we will
maintain and update the database to support quality of life features, align
with the models' input variables in case they end up changing, and add data
from the ongoing NFL season.

\section{Conclusion}

Regardless of the setting, whether it be trading, education, or coaching, the
models' estimates should not override informed decision-making or replace
cognitive reasoning, but instead assist in them. Thus, the way in which the
tool is presented and curated should encourage users to think about how the win
probabilities should be considered within the context of the game state they
correspond to. This philosophy will inform our approach to designing the
website. If this approach is successful, our tool should be ready for
responsible and informed usage for multiple purposes by a wide range of people
among the millions that are interested in American football.

\end{document}
